\section{Localized propagators and coarse covariances}\label{sec:fermion-primitives}

We use the local integration construction of \cite[Section~11]{QED3Manuscript}. The walk expansion provides its localization. A parameter is attached to each $M$-cube traversed
by a walk, together with its one-layer neighbourhood. Setting parameters
to zero then separates the operator into components. The coarse
covariance is obtained from this expansion by an exact resolvent
identity. At parameter value one it agrees with the exact covariance;
its local terms expand the determinant and the interaction.
We work at mesh $\eta_k=L^{-k}$ on the torus of side $\ell$; this is
the side denoted by $R$ in the local inverse construction.
Throughout this section the mass satisfies $0<\mu\le1$,
the range of the renormalization data.
Transposition is with respect to the bilinear lattice pairings:
for $O:E\to E'$, $\langle v,Ou\rangle_{E'}=
\langle O^Tv,u\rangle_E$. These pairings use the counting,
chamber, or coarse lattice measure of the respective space.
No complex conjugation enters this operation.

\subsection{Weakened walks and polymer propagators}\label{subsec:weakened-walks}

The resolvent parameter ranges over $Y_0=\R$ at level zero and
$Y_k=\R\cup\{\infty\}$ at $k\ge1$, with $-\infty=\infty$
as in the preceding resolvent construction.

Let $\pi_k$ be the partition into $M$-cubes. For integers $j\ge0$, write
$\square^{\sim j}$ for the union of partition cubes whose centres
are within periodic sup distance $jM$ of the centre of $\square$.
When this enlargement does not wrap, it has side $(2j+1)M$.
Let $\mathcal P_k$ be the set of
nonempty unions of these cubes connected through faces. For $X\in\mathcal P_k$
write $|X|_M$ for its number of cubes. With the net and length
conventions of \cref{app:polymer-geometry}, set
\[
 M d_M(X)=\inf\{\operatorname{length}(\mathcal N):
       \mathcal N\text{ is a net meeting every closed cube of }X\}.
\]
For a face-connected union of cubes, joining successive cube centres gives
$d_M(X)\le |X|_M-1$.

For a walk $\omega=(\square_0,\ldots,\square_n)$ from
\cref{subsec:propagator-walks}, put
\[
 X_\omega=\bigcup_{i=0}^n\square_i^{\sim1},\qquad
 s_\omega=\prod_{\square\subset X_\omega}s_\square,
 \qquad |s_\square|\le M^{\alpha_0},\quad \alpha_0=1/54.
\]
The union $X_\omega$ is a polymer and
\begin{equation}\label{eq:walk-polymer-size}
 |X_\omega|_M\le27(n+1),\qquad d_M(X_\omega)\le27(n+1).
\end{equation}
Indeed, each one-layer neighbourhood contains $27$ cubes. Successive
centres differ by at most $M$ in each coordinate, so their neighbourhoods
are joined through faces. The preceding bound on $d_M$ applies.
Every walk, including a walk of length zero, carries the parameters of
all cubes in $X_\omega$.

The left parametrix gives terms
\[
 S'_{y,\omega}=(h_nS_n)(-C_nh_{n-1}S_{n-1})\cdots
                 (-C_1h_0S_0)h_0.
\]
Here, as for $S_{y,\omega}$, $S_i$ is the local inverse on
$\square_i$ and $C_i=-[F_y,h_i]$. Define
\begin{align}
 S_y(s,A)&=\sum_\omega s_\omega S_{y,\omega}(A),&
 S_y(X,A)&=\sum_{X_\omega=X}S_{y,\omega}(A),\label{eq:weakened-walk-definitions}\\
 S'_y(s,A)&=\sum_\omega s_\omega S'_{y,\omega}(A),&
 S'_y(X,A)&=\sum_{X_\omega=X}S'_{y,\omega}(A).\notag
\end{align}
Their averages are denoted by $S_y^{\mathrm{sy}}$. The superscript
$\#$ below means any of these three choices. Put
$\kappa'=\kappa_{\mathrm w}+\kappa_0$, where
$\kappa_0>0$ and $\kappa_{\mathrm w}=\kappa+\kappa_0+3$, with $\kappa>0$, are fixed polymer decay rates.
Choose $M$ large enough that
\begin{equation}\label{eq:weakened-walk-contraction}
 q(M^{27\alpha_0})\le\tfrac12,\qquad
 q(e^{27\kappa'})\le\tfrac12,
 \qquad q(p)=54\pi C_Va'p/M.
\end{equation}

\begin{proposition}\label{prop:weakened-polymer-propagators}
For $A\in\widetilde{\mathcal G}_k$, $y\in Y_k$, and
$s\in\prod_{\pi_k}\{|s_\square|\le M^{\alpha_0}\}$, the series
\eqref{eq:weakened-walk-definitions} converge absolutely in $N_\gamma$,
uniformly in $(s,A)$, and
\begin{align}
 N_\gamma(S_y^\#(s,A))&\le32a'M^{27\alpha_0}w_k(y),
                                       \label{eq:weakened-propagator-bound}\\
 N_\gamma(\delta_{\alpha,A}S_y^\#(s,A))&\le304a'M^{27\alpha_0},\notag\\
 \|\mathbf1_{\Delta_z}S_y^\#(X,A)\mathbf1_{\Delta_{z'}}\|
   &\le32a'e^{27\kappa'}w_k(y)
       e^{-\gamma d(z,z')-\kappa'd_M(X)}.
                                       \label{eq:polymer-propagator-bound}
\end{align}
The norm in the last line is the sup operator norm. At $k=0$ it can
also be replaced by the supremum of the normalized kernel over the two
clusters, with $w_0(y)=(1+|y|)^{-1}$.
At parameter value one,
\begin{equation}\label{eq:polymer-propagator-sum}
 S_y^\#(1,A)=S_y(A)=\sum_{X\in\mathcal P_k}S_y^\#(X,A).
\end{equation}
The weakened operators are continuous on the closed parameter
polydisk, jointly analytic in its interior and in $A$, and analytic in
$A$ for each fixed parameter value, including boundary values.
For finite real $y$ the families are continuous in $y$.

The operator $S_y^\#(X,A)$ is supported on $X$ in both arguments and
depends on $A$ only on bonds with endpoints in $X$. If parameters vanish
off a union $Y$ of $M$-cubes, then each component of $Y$ is preserved by
$S_y^\#(s,A)$. On such a component the operator depends only on its
parameters and its connection. Its row at $x=(c,g)$ vanishes unless all
$M$-cubes meeting $K_c=\{c':|c'-c|_\infty\le2L\}$ have active
parameters; the corresponding assertion holds for columns.
\end{proposition}

\begin{proof}
By \eqref{eq:walk-polymer-size}, $|s_\omega|\le p_s^{n+1}$ with
$p_s=M^{27\alpha_0}$. Apply the walk row estimate with this weight.
The length-$n$ contribution is bounded by
$16a'p_sw_k(y)2^{-n}$; its H\"older row estimate is bounded by
$152a'p_s2^{-n}$. Summing these bounds proves the first two estimates.
For a fixed polymer use
\[
 \mathbf1[X_\omega=X]\le
       e^{-\kappa'd_M(X)}p_X^{n+1},\qquad p_X=e^{27\kappa'}.
\]
The same row estimate with $p_X$ proves
\eqref{eq:polymer-propagator-bound}. The left walk estimates follow
from the left parametrix by the same weighted product estimate; taking
the average preserves all bounds.

On a fixed torus, summing the row bounds over its finitely many unit
blocks also bounds the sum of the $N_\gamma$ norms of the individual
terms. The tails decay geometrically, uniformly in $(s,A)$. The series
may therefore be regrouped by $X_\omega$. At $s=1$ they are the right
and left inverse expansions, so both equal $S_y(A)$. This proves
\eqref{eq:polymer-propagator-sum}. Each finite sum is analytic in its
connection variables and polynomial in $s$. The uniform tail estimates
give the stated analyticity and continuity.
The local inverse is continuous at every real resolvent parameter.
Its walk terms inherit this property, and their tails are uniform in
$y$ since $w_k(y)\le1$. The sums are therefore continuous in real $y$.

A walk term is a product of local factors on the neighbourhoods
$\square_i^{\sim1}$. Its row and column cutoffs, its commutators, and
all connection variables consequently lie in $X_\omega$. If the first
row is nonzero, its photon site satisfies
$|c-z_{\square_0}|_\infty<3M/4$. Since $3M/4+2L<3M/2$, every
$M$-cube meeting $K_c$ is in $\square_0^{\sim1}$, and its parameter
occurs in $s_\omega$. The last cutoff gives the column assertion;
the order is reversed for a left walk. Finally, a surviving walk has
every cube of $X_\omega$ in $Y$. Since $X_\omega$ is face-connected,
it lies in a single component of $Y$. Its local factors use only the
connection in that component. This proves all the locality statements.
\end{proof}
The lattice symmetries used in this section have coordinate sign
changes and translations in $M\Z^3$. Their spin lifts lie in the
group generated by $\gamma_1,\gamma_2,\gamma_3,\gamma_5$; they act
on $(c,g)$ by $(rc+a,rg)$, where
$S\gamma_\mu S^{-1}=r_{\mu\mu}\gamma_\mu$ determines $r$.
Write $\widetilde g=(S,a)$ for such a symmetry.
The step group restricts $a$ to
$LM\Z^3$. A lift commuting with $\gamma_5$ preserves $y$; a lift
anticommuting with it sends $y$ to $-y$ at nonneutral values.
Gauge covariance follows term by term: the conjugating phases of
adjacent factors cancel. A lattice symmetry permutes cubes and walks,
and hence acts on the parameters by
$(\widetilde g s)_\square=s_{\widetilde g^{-1}\square}$.
Transposition reverses the order of
the local factors and interchanges the right and left expansions.
More precisely, with $\mathsf C=i\gamma_2$ one has
$F_y(A)^T=\mathsf C^{-1}F_{-y}(-A)\mathsf C$. Therefore
\begin{equation}\label{eq:weakened-transpose}
 S_y(s,A)^T=\mathsf C^{-1}S'_{-y}(s,-A)\mathsf C.
\end{equation}
Thus at the neutral values $y=0,\infty$ for $k\ge1$, and $y=0$
for $k=0$, the symmetrized family has the same transpose covariance as
the exact inverse. This is the reason for using it in the coarse
covariance and determinant.

\subsection{The fine and coarse resolvents}\label{subsec:fine-coarse-resolvent}

For $k\ge1$ let
\[
 \mathcal Q_k=Q_k(A)J^*,\qquad
 \mathcal Q_k^\sharp=JQ_k^T(-A),\qquad
 P_L=Q^T(-A)Q(A).
\]
Here $Q_k$ averages the fine lattice onto the unit lattice, and $Q$
averages the unit lattice onto the lattice of spacing $L$.
The notation $\sharp$ denotes the displayed algebraic transpose with
reversed connection; for complex $A$ it is not a Hermitian adjoint.
Here $B(z')=E_{z'}$ is the $L$-block from the preceding section.
To specify the phases, if $c\in\Delta_z$ and $z\in B(z')$, set
\[
 \varphi_A(z,c)=e^{ie_k\eta_k(\tau_kA)(z,c)},\qquad
 U_A(z',z)=e^{ie_k\eta_k(\tau A)(z',z)}.
\]
The path average $\tau A$ is the arithmetic mean of the oriented bond
sums in the six coordinate orders, and $\tau_kA$ is their sum along
the nested block centres. Then
\begin{align*}
 (\mathcal Q_kf)(z)&=\eta_k^3\sum_{c\in\Delta_z}
                \varphi_A(z,c)\frac1{48}\sum_gf(c,g),\\
 (\mathcal Q_k^\sharp v)(c,g)&=\varphi_A(z,c)^{-1}v(z),\\
 (P_Lv)(z)&=U_A(z',z)^{-1}L^{-3}
                    \sum_{z''\in B(z')}U_A(z',z'')v(z'').
\end{align*}
The nesting of the odd blocks and cancellation of the phases give
\begin{equation}\label{eq:block-algebra}
 \mathcal Q_k\mathcal Q_k^\sharp=1,\quad P_L^2=P_L,\quad
 \mathcal Q_k^\sharp\mathcal Q_k=\Pi_1(A),\quad
 \mathcal Q_k^\sharp P_L\mathcal Q_k=\Pi_L(A).
\end{equation}
All these operators commute with constant spin matrices.

Write $S_k=S_\infty(A)$ and define maps on the unit spinors by
\begin{align}
 H_k&=b_kS_k\mathcal Q_k^\sharp,&
 \mathfrak H_k&=b_k\mathcal Q_kS_k,\notag\\
 D_k&=b_k-b_k^2\mathcal Q_kS_k\mathcal Q_k^\sharp,&
 M_k&=b_k-D_k.\label{eq:coarse-dirac-definition}
\end{align}
For real $y$ put
\[
 \mathcal B_{k,y}=(b_k+i\gamma_5y)^{-1}(1-P_L)
                   +(b_k+b_L+i\gamma_5y)^{-1}P_L.
\]
These $4\times4$ inverses exist since
$(t+i\gamma_5y)^{-1}=(t-i\gamma_5y)/(t^2+y^2)$ for $t>0$.
The idempotent decomposition and \eqref{eq:block-algebra} give
\begin{equation}\label{eq:fine-resolvent-perturbation}
 \mathcal B_{k,y}^{-1}=b_k+b_LP_L+i\gamma_5y,\qquad
 F_y=F_\infty-b_k^2\mathcal Q_k^\sharp
                       \mathcal B_{k,y}\mathcal Q_k.
\end{equation}

\begin{proposition}\label{prop:fine-coarse-resolvents}
For $k\ge1$, $A\in\widetilde{\mathcal G}_k$, and real $y$, the
operator $D_k+b_LP_L+i\gamma_5y$ is invertible and its inverse is
\begin{equation}\label{eq:coarse-resolvent-formula}
 \Gamma_{k,y}=\mathcal B_{k,y}
       +b_k^2\mathcal B_{k,y}\mathcal Q_kS_y
                         \mathcal Q_k^\sharp\mathcal B_{k,y}.
\end{equation}
Moreover,
\begin{equation}\label{eq:resolved-fine-covariance}
 S_y=S_k+H_k\Gamma_{k,y}\mathfrak H_k.
\end{equation}
In particular, $\Gamma_k=\Gamma_{k,0}$ exists and
$\delta Z_k(A)=\det(D_k+b_LP_L)\ne0$.
\end{proposition}

\begin{proof}
Set $Z=b_k^2\mathcal Q_k^\sharp\mathcal B_{k,y}\mathcal Q_k$.
The two fine inverses exist by the preceding section, and
\eqref{eq:fine-resolvent-perturbation} gives both identities
\[
 S_y-S_k=S_kZS_y=S_yZS_k.
\]
Since $D_k+b_LP_L+i\gamma_5y
=\mathcal B_{k,y}^{-1}-b_k^2\mathcal Q_kS_k\mathcal Q_k^\sharp$,
multiplication by the right side of
\eqref{eq:coarse-resolvent-formula} gives the identity: the additional
term is
$b_k^2\mathcal Q_k(S_y-S_k-S_kZS_y)
\mathcal Q_k^\sharp\mathcal B_{k,y}=0$.
The product in the other order uses $S_y-S_k-S_yZS_k=0$.
This proves the inverse formula without assuming that the coarse
operator is invertible. Substituting the formula in
$H_k\Gamma_{k,y}\mathfrak H_k$ gives
\[
 S_kZS_k+S_kZS_yZS_k
     =S_kZ(S_k+S_yZS_k)=S_kZS_y=S_y-S_k.
\]
At $y=0$ the coarse operator is an invertible map on the finite
dimensional unit spinor space, so its determinant is nonzero.
\end{proof}

At level zero this algebra is replaced by a single chamber
integration. The operators $b_0$, $S_\infty$, $\mathcal B_{0,y}$, and
$M_0$ are not defined. Instead,
\begin{equation}\label{eq:level-zero-covariance}
 F_0=D_A+\mu+b_LJP_LJ^*,\qquad
 \Gamma_0=F_0^{-1},\qquad \delta Z_0(A)=\det F_0\ne0.
\end{equation}
Here $\Gamma_0$ acts on $192$ components per photon site, whereas
$\Gamma_k$ for $k\ge1$ acts on four components per unit site.

\subsection{Localized coarse covariances}\label{subsec:localized-coarse-covariance}

The weakened families for $k\ge1$ are
\begin{align}
 H_k(s,A)&=b_kS_\infty(s,A)\mathcal Q_k^\sharp,
                                          \label{eq:local-coarse-families}\\
 \Gamma_k(s,A)&=\mathcal B_{k,0}
    +b_k^2\mathcal B_{k,0}\mathcal Q_kS_0^{\mathrm{sy}}(s,A)
                 \mathcal Q_k^\sharp\mathcal B_{k,0},\notag\\
 M_k(X,A)&=b_k^2\mathcal Q_kS_\infty^{\mathrm{sy}}(X,A)
                  \mathcal Q_k^\sharp.\notag
\end{align}
The unprimed polymer propagators $S_{k,y}(X,A)=S_y(X,A)$ are used
in the determinant integral. At level zero let
$\operatorname{Cub}(c)$ be the $M$-cubes having a lattice site in
common with $K_c$, and set
\begin{equation}\label{eq:local-level-zero-families}
 \begin{aligned}
 (H_0(s,A)u)(c,g)&=
       \left(\prod_{\square\in\operatorname{Cub}(c)}s_\square\right)u(c,g),\\
 \Gamma_0(s,A)&=S_0^{\mathrm{sy}}(s,A),\qquad
 S_{0,y}(X,A)=S_y(X,A).
 \end{aligned}
\end{equation}
At $s=1$ these are the exact operators; also
$\sum_XM_k(X,A)=M_k(A)$. The exponential formulas for the averaging
phases and \cref{prop:weakened-polymer-propagators} give
analyticity on the same domains.

For use in the estimates, let
$\mathfrak i(A)=e_k\|\operatorname{Im}A\|_\infty$.
The small-connection domain and the charge bound give
$\mathfrak i(A)\le1/(6L)$; reduce $\gamma$, if necessary, so that
$\gamma L\le1/6$. The path lengths imply
\begin{equation}\label{eq:coarse-phase-bounds}
 |\varphi_A(z,c)^{\pm1}|\le e^{3\mathfrak i(A)/2},\qquad
 \|\mathcal Q_k\|_{\infty\to\infty},
 \|\mathcal Q_k^\sharp\|_{\infty\to\infty}
       \le e^{3\mathfrak i(A)/2}.
\end{equation}
The sup norms here are formed with the Euclidean spinor norm.
Each unit block has chamber measure one. In particular, for a fine
operator $O$,
\begin{equation}\label{eq:block-compression-bound}
 |(\mathcal Q_kO\mathcal Q_k^\sharp)(z,z')|_2
       \le \exp(1)\,\|\mathbf1_{\Delta_z}O\mathbf1_{\Delta_{z'}}\|.
\end{equation}
Indeed, apply the two averaging bounds to a unit spinor supported at
$z'$; their product is $e^{3\mathfrak i(A)}\le \exp(1)$.

\begin{proposition}\label{prop:coarse-local-family-bounds}
The families \eqref{eq:local-coarse-families} satisfy
\begin{align}
 \|H_k(s,A)\|_{\infty\to\infty}
       &\le32\exp(1)b_ka'M^{27\alpha_0},\label{eq:weakened-H-bound}\\
 \|\delta_{\alpha,A}H_k(s,A)\|_{\infty\to\infty}
       &\le304\exp(1)b_ka'M^{27\alpha_0},\notag\\
 |\Gamma_k(s,A)(z,z')|_2
       &\le C_\Gamma M^{27\alpha_0}e^{-\gamma d(z,z')},
                                             \label{eq:weakened-Gamma-bound}\\
 |M_k(X,A)(z,z')|_2
       &\le32\exp(1)b_k^2a'e^{27\kappa'}
                     e^{-\gamma d(z,z')-\kappa'd_M(X)},\notag
\end{align}
where
\[
 C_\Gamma=\frac{2(1+2\exp(1))}{b_-}+128\exp(1)(1+2\exp(1))^2a'.
\]
The H\"older estimate is on short pairs, including the distinct
sites of one cluster. The families $H_k(s,A)$ and $\Gamma_k(s,A)$
preserve every component of the active $M$-cubes and depend there only
on its parameters and connection. The row of $H_k(s,A)$ at $(c,g)$
vanishes unless all cubes meeting $K_c$ are active. The polymer
$M_k(X,A)$ is supported on the unit sites of $X$ and depends only on
the connection in $X$.

At level zero the normalized kernel of $\Gamma_0(s,A)$ is bounded by
$32a'M^{27\alpha_0}e^{-\gamma d(c,c')}$, and
\eqref{eq:polymer-propagator-bound} bounds $S_{0,y}(X,A)$ with its
factor $(1+|y|)^{-1}$. The multiplier $H_0(s,A)$ has sup norm
at most $M^{8\alpha_0}$ and equals the identity at $s=1$.
Its row at $(c,g)$ vanishes unless every cube meeting $K_c$
is active. There are no short pairs of distance less than one
at this level.
\end{proposition}

\begin{proof}
The first two estimates follow by composing the weakened fine bounds
with \eqref{eq:coarse-phase-bounds}. The third and fourth use
\eqref{eq:block-compression-bound}. For the covariance it remains to
control the two factors $\mathcal B_{k,0}$. Its kernel is
\[
 \mathcal B_{k,0}(z,z')=b_k^{-1}\mathbf1[z=z']1_4
      -\frac{b_L}{b_k(b_k+b_L)}P_L(z,z').
\]
It vanishes unless $z,z'$ belong to one $L$-block. Since
$|P_L(z,z')|_2\le\exp(1)L^{-3}$ on that block, its row and column sums
are at most $(1+2\exp(1))/b_k$. In a nonzero term of
$\mathcal B_{k,0}G\mathcal B_{k,0}$,
where $G=\mathcal Q_kS_0^{\mathrm{sy}}(s,A)\mathcal Q_k^\sharp$,
the two intermediate sites are within $L$ of the exterior sites.
Therefore
\[
 e^{-\gamma d(z_1,z_2)}\le
      e^{2\gamma L}e^{-\gamma d(z,z')}
      \le4e^{-\gamma d(z,z')}.
\]
The fine bound and \eqref{eq:block-compression-bound} give
$|G(z_1,z_2)|_2\le32\exp(1)a'M^{27\alpha_0}e^{-\gamma d(z_1,z_2)}$.
Multiplying the row and column sums cancels the factor $b_k^2$ and
gives $128\exp(1)(1+2\exp(1))^2a'M^{27\alpha_0}$.
The local term $\mathcal B_{k,0}$ is bounded by
$2(1+2\exp(1))b_-^{-1}e^{-\gamma d(z,z')}$, proving
\eqref{eq:weakened-Gamma-bound}.

The odd block partitions are nested, and every unit block and every
$L$-block belongs to a single $M$-cube. Moreover, the paths entering
$\varphi_A(z,c)$ lie in $\Delta_z$, and the paths entering
$U_A(z',z)$ lie in the fine block of side $L$ centred at $z'$:
each coordinate of a rectilinear path remains between those of its
endpoints. Consequently the averaging maps and $\mathcal B_{k,0}$
preserve every union of $M$-cubes and depend only on its connection.
Composing them with the component-local fine families proves the
locality assertions. This applies also to the parameter-independent
term $\mathcal B_{k,0}$ in $\Gamma_k(s,A)$.
For the row of $H_k$, the terminal average does not change the first
row index of $S_\infty(s,A)$, so the stronger vanishing statement is
preserved. At level zero the assertions for $\Gamma_0$ follow
from the fine kernel bound and locality with $\eta_0=1$.

For $H_0$, lift $c$ and $K_c$ to the lattice. In each coordinate
$K_c$ is an interval of $4L+1$ sites. Since $M\ge L^2$ and $L\ge5$,
this interval meets at most two consecutive $M$-blocks. The cubes
meeting $K_c$ are the product of these three sets of blocks.
They number at most eight, form a face-connected set, and include
the cube of $c$. Thus the defining product is a monomial of degree
at most eight, bounded by $M^{8\alpha_0}$. If one of its cubes is
inactive, its factor is zero. If all are active, their connectedness
puts them in the active component containing $c$; hence changing
parameters outside that component has no effect on the row or
column at $(c,g)$. The multiplier is scalar in spin and chamber
labels and independent of the connection. These observations give
its analyticity, gauge covariance, and component locality.
\end{proof}

These families have the gauge and lattice covariance of their exact
counterparts. For $\Gamma_k$ and $M_k$, transposition follows from
\eqref{eq:weakened-transpose} at a neutral value and the scalar
averaging phases. The barred field map is
$\bar H_k(s,A)=\mathsf C^{-1}H_k(s,-A)\mathsf C$.
Its bounds follow from those for $H_k$ because $\mathsf C$ is unitary
and $\mathsf C\gamma_\mu\mathsf C^{-1}=-\gamma_\mu^T$.
When component sup norms are used in the Grassmann substitution,
the Euclidean operator bounds above acquire at most the factor two
from $|v|_c\le|v|_2\le2|v|_c$.

